% !TeX root = ../main.tex
\section{Разделяющиеся переменные}

\subsection{Дифференциальные уравнения с разделяющимися переменными}

Уравнение вида

\[
f_1(x)g_1(y) + f_2(x)g_2(y) = 0
\]

Для решения \textit{разделяем} части, содержащие $x$ и $y$, в разные стороны, затем интегрируем.

Например [C2 | 4]:

\[
y'\cos{x}+y(1+y)\sin{x}=0
\]

$y=0$ или $y=-1$ являются решениями, далее рассматриваем иные случаи:

\[
\frac{dy}{y(y+1)}=-\tan{x}dx
\]

Интегрируем обе части:

\[
\ln{|y|}-\ln{|y+1|}=\ln{|\cos{x}|}+C
\]

Потенцируем, то есть применяем $\exp$ к обеим частям:

\[
\left|\frac{y}{y+1}\right|=e^C|\cos{x}|
\]

Заметим что можно убрать модули за счет константы:

\[
\frac{y}{y+1}=C\cos{x} \Rightarrow y=c(y+1)\cos{x}
\]

Это и будет решение в неявном виде (в ответ также нужно будет добавить $y=0$ и $y=-1$).

\subsection{Разделение переменных линейной заменой аргумента}

Уравнение вида

\[
\frac{dy}{dx} = f(ax + by)
\]

\textbf{Замена} $t = ax + by$, после которой уравнение приводится к виду с разделяющимися переменными:

\[
dy=f(t)dx, dt=adx+bdy \Rightarrow dt=adx+bf(t)dx \Rightarrow \frac{dt}{a+bf(t)}=dx
\]

Например [Ф | 62]:

\[
y'=\cos(y-x)
\]

Замена $t=y-x \Rightarrow dt=dy-dx$

\[
dy=\cos{t}dx, dt=dy-dx \Rightarrow dt=\cos{t}dx-dx \Rightarrow dt=(\cos{t}-1)dx
\]

$\cos{t}-1=0 \Leftrightarrow t=\pi + 2\pi k, k \in \mathbb{Z}$ -- решение, далее не рассматриваем.

Далее нам понадобится тождество $\cos{t}-1=-2\sin^2{\frac{t}{2}}$.

\[
\frac{dt}{\cos{t}-1}=dx \Rightarrow \frac{dt}{-2\sin^2{\frac{t}{2}}}=dx \Rightarrow \ctg{\frac{t}{2}}=x+C
\]

Итого ответ: $
\begin{cases} 
    \ctg{\frac{y-x}{2}}=x+C \\ 
    y=2\pi k, k\in \mathbb{Z}
\end{cases}
$

\section{Однородные уравнения}

\subsection{Однородное уравнение}

Уравнение вида
\[
M(x,y)dx + N(x,y)dy = 0,
\]

где $M$ и $N$ непрерывные в некоторой области $D \in \mathbb{R}^2$ однородные функции одинаковой степени. С помощью \textbf{замены} $y=x\cdot U(x)$ сводится к уравнению с разделяющимися переменными $x$ и $U$.

\textit{В задачах из задавальника чаще всего однородное уравнение уже имеет вид уравнения с разделяющимися переменными и нужно лишь для решения \textit{линейного уравнения}.}

\section{Линейные уравнения и приводящиеся к ним.}

\subsection{Линейные уравнения первого порядка.}

Уравнение вида

\[
y'+P(x)y=Q(x).
\]

Чаще всего решается методом \textit{вариации постоянной}: сначала ищем решение однородного уравнения $y'+P(x)y=0$, получая некое выражение с константой. Эту констатну $C$ варируют, то есть считают функцией от $x$. Далее находят эту функцию с помощью подстановки в исходное уравнение.

Например [C3 | 28]:

\[
xy'+(1+x^2)y+x=0
\]

\[
y'+\frac{1+x^2}{x}\cdot y=-1 \text{-- линейное уравнение (л.у.)}
\]

Рассматриваем соответствующее однородное уравнение:

\[
y'+\frac{1+x^2}{x}\cdot y=0 \text{-- р.п.}
\]

\[
\frac{dy}{y}=-\frac{1+x^2}{x}dx\ \Rightarrow \ \ln{|y|}=-\ln{|x|}-\frac{x^2}{2} \ \Rightarrow \ y=\frac{c}{x}\cdot e^{-\frac{x^2}{2}} 
\]

Теперь варируем постоянную: $y=\frac{c(x)}{x}\cdot e^{-\frac{x^2}{2}}$ и подставляем в исходное уравнение чтобы найти $c(x)$. Для этого вычислим $y'(x)$ из последнего выражения:

\[
y'=\frac{c'(x)}{x}\cdot e^{-\frac{x^2}{2}}-\frac{c(x)}{x^2}\cdot e^{-\frac{x^2}{2}}-c(x)\cdot e^{-\frac{x^2}{2}}
\]

Подставляем в исходное, и сократив получаем:

\[
\frac{c'(x)}{x}e^{-\frac{x^2}{2}}=-1 \ \Rightarrow \ c'(x)=-xe^{\frac{x^2}{2}} \ \Rightarrow \ c(x)=-e^{\frac{x^2}{2}}+c_0
\]

Ну и отсюда получаем
\[
y=\frac{(-e^{\frac{x^2}{2}}+c_0)\cdot e^{-\frac{x^2}{2}}}{x}=\frac{c_0e^{\frac{x^2}{2}}-1}{x}
\]

Ответ: $y=\frac{c_0e^{\frac{x^2}{2}}-1}{x}$

\subsection{Уравнение Бернулли}

Уравнение вида

\[
y'+P(x)y=Q(x)y^m
\]

Приводится к линейному с помощью \textbf{замены} $z(x)=\frac{1}{y^{m-1}}$ (фактически делим на $y^m$ чтобы справа была функция от $x$ а потом заменяем на $z$ то что стоит рядом с $P(x)$).

\subsection{Уравнение Риккати}

Уравнение вида

\[
y'+a(x)y+b(x)+c(x)y^2=0
\]

Ищем частное решение $y_0$ (можно искать подбором, например $x^m$, $e^{\alpha x}$, $ax+b$, $a\sin{kx}+b\cos{kx}$, ...). \textbf{Замена:} $y=y_0+z(x)$ сводит уравнение к уравнению Бернулли.

Например [C3 | 90]:

\[
y'+e^{-x}y^2+y=3e^x
\]

Заметим что $y_0=e^x$ является решением. Тогда замена $y=e^x+z(x)$. $y'=e^x+z'$

\[
e^x+z'+e^{-x}(e^x+z)^2+e^x+z=3e^x
\]

\[
z'+3z=e^{-x}z^2 \text{ - уравнение Бернулли}
\]

\[
\frac{z'}{z^2}+\frac{3}{z}=e^{-x}
\]

\[
t=\frac{1}{z}\ \Rightarrow \ t'-3t=-e^{-x} \text{ - л.у.} 
\]

Однородное:
\[
t'-3t=0\ \Rightarrow \ t=ce^{3x}
\]

Варируем постоянную:

\[
t=c(x)e^{3x} \ \Rightarrow \ t'=c'e^{3x} + 3ce^{3x}
\]

\[
-c'e^3x-3ce^{3x}+3ce^{3x}=e^{-x} \ \Rightarrow \ c'=-e^{-4x} \ \Rightarrow \ c(x)=\frac{1}{4}\cdot e^{-4x}+c_0
\]

\[
t=\left(\frac{1}{4}\cdot e^{-4x}+c_0\right)e^{3x}=\frac{1}{4}\cdot e^{-x}+c_0e^{3x}
\]

\[
z=\frac{4}{e^{-x}+ce^{3x}}
\]

Ответ: $
\begin{cases} 
    y=e^x+\frac{4}{e^{-x}+ce^{3x}} \\ 
    y=e^x
\end{cases}
$

\section{Уравнение в полных дифференциалах. Интегрирующий множитель}

\subsection{Уравнение в полных дифференциалах}

Уравнение вида
\begin{equation}
    M(x, y)dx+N(x, y)dy=0
    \label{eq:pd}
\end{equation}

Если его левая часть является полным дифференциалом некоторой функции, то есть

\[
M(x, y)dx+N(x, y)dy=d\Phi(x, y)
\]

Если функции $M, N, \frac{\partial N}{\partial x}, \frac{\partial M}{\partial y}$ непрерывны в некоторой односвязной области $D\subset \mathbb{R}^2$, то условие

\[
\frac{\partial N}{\partial x}=\frac{\partial M}{\partial y}
\]

является необходимым и достаточным для того чтобы уравнение было уравнением в полных дифференциалах.

\subsection{Интегрирующий множитель}

Функция $\mu=\mu(x, y)\not = 0$, после умножения на которую уравнение вида \eqref{eq:pd} превращается в уравнение в полных дифференциалах.

\thm{}{
  Если функции $M$ и $N$ непрерывны, имеют непрерывные частные производные и $M^2+N^2\not =0$, то интегрирующий множитель существует.
}{}

